% =====================================================================
% uni-thesis-engine -- modular, topic-agnostic thesis template.
% Compiles with pdflatex or xelatex using only standard TeX Live fonts.
% Every result figure, number and citation is injected from generated/:
%   generated/auto_metrics.tex     result macros (\Metric..., \Stat..., \Diff...) + flags
%   generated/auto_tables.tex      booktabs three-line tables
%   generated/auto_figures.tex     fig_foundry vector figures
%   generated/auto_literature.tex  DOI-verified related work
% Never type an experimental number into chapters/ -- the audit will reject it.
% =====================================================================
\documentclass[11pt,a4paper,oneside]{report}

\usepackage{iftex}
\ifPDFTeX
  \usepackage[T1]{fontenc}
  \usepackage[utf8]{inputenc}
  \usepackage{lmodern}
\else
  \usepackage{fontspec}
\fi
\usepackage[margin=2.5cm]{geometry}
\usepackage{amsmath,amssymb}
\usepackage{booktabs}
\usepackage{graphicx}
\usepackage[round,authoryear]{natbib}
\usepackage[hidelinks]{hyperref}

\input{generated/auto_metrics.tex}

\title{\ThesisTitle}
\author{Author Name\\ \small Supervisor: Supervisor Name}
\date{\today}

\begin{document}
\maketitle

\begin{abstract}
This thesis investigates \ThesisTitle. A reproducible pipeline retrieves and DOI-verifies the
related literature, executes every experiment from a declarative topic profile, routes each
comparison to the statistical test its assumptions allow, and injects all reported values into
this document through auto-generated macros. On the primary metric (\MetricPrimaryName), the
proposed configuration (\MetricPrimaryProposedName) reaches \MetricPrimaryProposedMean{} $\pm$
\MetricPrimaryProposedStd{} \MetricPrimaryUnit, compared with \MetricPrimaryRunnerUpMean{} for the
strongest alternative (\MetricPrimaryRunnerUpName; $p$ \StatPrimaryPval, \StatPrimaryTest).
\end{abstract}

\tableofcontents

\chapter{Introduction}
\label{chap:intro}

\section{Background and motivation}
This thesis addresses the topic \emph{\ThesisTitle}. Reliable empirical claims require that every
reported value can be traced to an executed experiment and every cited work to a verifiable
publication. The study is therefore organised around an automated, auditable research pipeline.

\section{Research objectives}
\begin{itemize}
  \item Establish a reproducible experimental protocol driven entirely by a declarative topic profile.
  \item Quantify the proposed configuration against competing alternatives on \MetricPrimaryName{}
        and the secondary metrics, with assumption-aware statistical testing.
  \item Report only values generated by the experiments, with provenance hashes linking the
        manuscript to the result bundle.
\end{itemize}

\section{Thesis structure}
Chapter~\ref{chap:related} reviews the DOI-verified literature. Chapter~\ref{chap:method} describes
the methodology and statistical protocol. Chapter~\ref{chap:experiments} reports the experimental
results, and Chapter~\ref{chap:conclusion} concludes the work.

\chapter{Related Work}
\label{chap:related}

The following works were retrieved from OpenAlex and Crossref and kept only after their DOI
metadata (title, authors and year) were confirmed on Crossref. Limitations are quoted verbatim from
the published abstracts when available; otherwise they await manual review.

\input{generated/auto_literature.tex}

\chapter{Methodology}
\label{chap:method}

\section{Experimental protocol}
All settings of the study (data source, competing methods, metrics, number of repetitions and
random seed) are declared in the topic profile \texttt{\ThesisProfileId}. Each configuration is
evaluated over independent repeated runs, and every run is recorded in a result bundle stamped
with the source-code commit and the SHA-256 digest of the input data.

\ifAutoHasRsm
\section{Response surface methodology}
The process factors are explored with a central composite design comprising factorial, axial and
replicated centre points. The response is modelled by the second-order polynomial
\begin{equation}
  \hat{y} = \beta_{0} + \sum_{i} \beta_{i} x_{i} + \sum_{i} \beta_{ii} x_{i}^{2}
            + \sum_{i<j} \beta_{ij} x_{i} x_{j},
\end{equation}
fitted by ordinary least squares. Model significance is assessed by an $F$-test, adequacy by the
lack-of-fit test against pure error, and the stationary point is classified from the eigenvalues
of the quadratic form.
\fi

\ifAutoHasAblation
\section{Ablation design}
The contribution of each core module is isolated by evaluating cumulative or leave-one-out
configurations of the full system under identical conditions.
\fi

\section{Statistical analysis}
For every metric, normality is examined with the Shapiro--Wilk test and homogeneity of variance
with Levene's test. Two-group comparisons are routed to the paired or Student $t$-test, Welch's
$t$-test, the Wilcoxon signed-rank test or the Mann--Whitney $U$ test according to these
diagnostics, and are accompanied by Cohen's $d$ or the rank-biserial correlation. Comparisons of
several groups use one-way ANOVA with Tukey's HSD or the Kruskal--Wallis $H$ test. Significance
markers follow the conventional thresholds stated beneath each table.

\chapter{Experiments and Results}
\label{chap:experiments}

\section{Overview}
This chapter reports the outcomes of the experiments declared in profile \texttt{\ThesisProfileId}.
Every value below is a macro generated from the signed result bundle; tables and figures are
produced by the same pipeline.

\input{generated/auto_tables.tex}

\ifAutoHasComparison
\section{Comparison with baseline methods}
Table~\ref{tab:auto-comparison} summarises all metrics. On \MetricPrimaryName, the proposed
method (\MetricPrimaryProposedName) achieves \MetricPrimaryProposedMean{} $\pm$
\MetricPrimaryProposedStd{} \MetricPrimaryUnit{} (confidence interval \MetricPrimaryProposedCiLow{} to
\MetricPrimaryProposedCiHigh), whereas the strongest baseline (\MetricPrimaryRunnerUpName) reaches
\MetricPrimaryRunnerUpMean{} $\pm$ \MetricPrimaryRunnerUpStd. The difference of \DiffPrimaryAbs{}
(\DiffPrimaryRelPct\,\% relative) is assessed with the \StatPrimaryTest{} ($p$ \StatPrimaryPval;
\StatPrimaryEffectName{} \StatPrimaryEffect; \StatPrimaryN{} runs per method). Across all methods,
the \StatPrimaryOmnibusTest{} yields $p$ \StatPrimaryOmnibusPval.
\fi

\ifAutoHasAblation
\section{Ablation study}
Table~\ref{tab:auto-ablation} reports the ablation configurations. The full system reaches
\MetricPrimaryProposedMean{} $\pm$ \MetricPrimaryProposedStd{} \MetricPrimaryUnit{} on
\MetricPrimaryName, against \MetricPrimaryRunnerUpMean{} for the best reduced configuration
(\MetricPrimaryRunnerUpName), a gain of \DiffPrimaryAbs{} (\DiffPrimaryRelPct\,\%; \StatPrimaryTest,
$p$ \StatPrimaryPval). The omnibus \StatPrimaryOmnibusTest{} across configurations gives
$p$ \StatPrimaryOmnibusPval.
\fi

\ifAutoHasRsm
\section{Response surface optimisation}
The design runs are listed in Table~\ref{tab:auto-design}. The fitted model
(Table~\ref{tab:auto-coef}) reads
\begin{equation}
  \StatRsmPrimaryEquation
\end{equation}
with \StatRsmPrimaryFactorLegend. The ANOVA in Table~\ref{tab:auto-anova} gives
$R^{2}$ = \StatRsmPrimaryRsq{} and adjusted $R^{2}$ = \StatRsmPrimaryRsqAdj{}
($F$ = \StatRsmPrimaryFvalue, $p$ \StatRsmPrimaryPval), and the lack-of-fit test returns
$p$ \StatRsmPrimaryLackOfFitPval. The stationary point is a \StatRsmPrimaryOptNature{} located at
\StatRsmPrimaryOptPoint, with a predicted response of \StatRsmPrimaryOptResponse{}
\MetricPrimaryUnit. Figure~\ref{fig:auto-rsm} visualises the surface.
\fi

\section{Figures}
\input{generated/auto_figures.tex}

\chapter{Conclusion}
\label{chap:conclusion}

\section{Summary}
This thesis studied \ThesisTitle{} with a fully traceable pipeline. On the primary metric
(\MetricPrimaryName), the proposed configuration (\MetricPrimaryProposedName) achieved
\MetricPrimaryProposedMean{} \MetricPrimaryUnit, improving on \MetricPrimaryRunnerUpName{} by
\DiffPrimaryAbs{} ($p$ \StatPrimaryPval).
\ifAutoHasRsm
The response-surface analysis located a \StatRsmPrimaryOptNature{} at \StatRsmPrimaryOptPoint{} with a
predicted \MetricPrimaryName{} of \StatRsmPrimaryOptResponse.
\fi

\section{Limitations and future work}
The conclusions are bounded by the data sources and operating ranges declared in the topic
profile. Future work should validate the findings on independent datasets, extend the factor
ranges where the optimum lies near a boundary, and examine additional metrics relevant to
deployment.


\bibliographystyle{plainnat}
\bibliography{refs}

\end{document}
